\documentclass{article}
\usepackage[utf8]{inputenc}
\usepackage[T1]{fontenc}
\usepackage{geometry}
\usepackage{graphicx}
\usepackage{float}
\usepackage{hyperref}
\usepackage[english,indonesian]{babel}

\title{Optimasi Pemrosesan Template LaTeX Menggunakan Pendekatan Berbasis AI}
\author{Academic Support Agent}
\date{\today}

\begin{document}
\maketitle

\begin{abstract}

\end{abstract}

\chapter{PENDAHULUAN}


\chapter{TINJAUAN PUSTAKA}


% Original content (commented out):
% % !TEX program = pdflatex
% !BIB program = biber
\input{setup}

\title{Partisanship and Its Effects on Inter-galatic Societies}
\author{Luke Skywalker}
\date{\today}

\doublespacing
\begin{document}
  
  \counterwithout{lstlisting}{chapter}
  \pagenumbering{roman}
  \include{frontmatter/titlepage}
  \afterpage{\null\addtocounter{page}{-1}\newpage}
  
  %\include{frontmatter/abstract}
  {\hypersetup{linkcolor=black}
    % This allows the paragraphs to be numbered (e.g. 1.2.3.4)
    \setcounter{secnumdepth}{4}
    
    % This allows the TOC to go down to subsection level, but not further
    \setcounter{tocdepth}{3}
    
    \tableofcontents
    \listoffigures
  }

  \setlength{\parindent}{1em}
  \setlength{\parskip}{0em}

  \newpage
  \pagenumbering{arabic}
  
  % You should always divide your documents into small parts and use \include to merge them together.
  \include{chapters/intro}
  
  \printbibliography

\end{document}


\end{document}